\section{Arithmetic profiles at every finite scale}\label{sec:arithmetic}
Now $\calA_\alpha=\{I,R_{2\pi\alpha_1},\ldots,R_{2\pi\alpha_r}\}$, with $r$ fixed. Define
\begin{equation}\label{eq:psi-eta}
 \psi_\alpha(H)=\min_{0<\norm{m}_1\le H}\norm{m\cdot\alpha}_{\R/\Z},
 \qquad
 \eta_\alpha(q)=\max_{1\le i\le r}\norm{q\alpha_i}_{\R/\Z}.
\end{equation}
These are finite functions of the alphabet parameters. The first may vanish in the presence of a rational relation. The upper profile remains usable in that case.

\begin{theorem}[Finite arithmetic bounds]\label{thm:arithmetic}
Let $\kappa=\rho/\sqrt2-2\varepsilon>0$. For a machine of peak $K$, put
$L=\lceil(2K)^{1/r}\rceil$ and $B=r(L-1)$. Then
\begin{equation}\label{eq:psi-lower}
 \lfloor N/B\rfloor\psi_\alpha(B)^2\le\log(1/\kappa).
\end{equation}
For every integer $q\ge4$,
\begin{equation}\label{eq:eta-upper}
 N\eta_\alpha(q)\le q^2/100
       \quad\Longrightarrow\quad W_{N,0}(\alpha)\le q.
\end{equation}
Consequently every finite planar rotation alphabet satisfies
\begin{equation}\label{eq:sqrt-upper}
                     W_{N,0}(\alpha)\le\max(4,\lceil\sqrt{50N}\rceil).
\end{equation}
All upper bounds use common command rows and hold also after adjoining reflection in the first coordinate axis.
\end{theorem}
\begin{proof}
Construct a word of length $B$ in $r$ consecutive segments of $L-1$ positions. In segment $i$ give $J_i$ copies of the $i$th rotation followed by idle commands, independently and uniformly for $J_i\in\{0,\ldots,L-1\}$. Differences of count vectors have $\ell_1$ norm at most $B$. If $\psi_\alpha(B)>0$, the $L^r\ge2K$ phases are distinct and separated by $\psi_\alpha(B)$. Corollary~\ref{cor:separated} and \eqref{eq:peak-profile}, with $-\log(1-x)\ge x$, give \eqref{eq:psi-lower}. If the minimum is zero the displayed inequality is automatic. All letters in a segment are actual commands, not a free count input.

For the upper bound let $P_q$ be the regular $q$-gon inscribed in the unit circle. Choose nearest integers $p_i$ to $q\alpha_i$, and put $\delta_i=2\pi(\alpha_i-p_i/q)$, so $|\delta_i|\le\pi/q$. Elementary chord geometry gives
\begin{equation}\label{eq:dilation}
 R_{2\pi\alpha_i}P_q\subseteq\Lambda_iP_q,
 \quad
 \Lambda_i=\frac{\cos(\pi/q-|\delta_i|)}{\cos(\pi/q)},
 \quad
 \log\Lambda_i\le\frac{4\pi^2\eta_\alpha(q)}{q^2}.
\end{equation}
Indeed the ray of a rotated vertex meets a neighboring edge at radial coordinate $\cos(\pi/q)/\cos(\pi/q-|\delta_i|)$. Integrating the derivative of $\log\cos(\pi/q-t)$ bounds the logarithm by $|\delta_i|\tan(\pi/q)$, and $\tan(\pi/q)\le2\pi/q$ for $q\ge4$.

Take $\Lambda=\max_i\Lambda_i$, including one for the identity. Under \eqref{eq:eta-upper}, $\Lambda^N\le e^{\pi^2/25}<2$. Since $P_q$ contains the half-unit disk, Theorem~\ref{thm:enclosure} with $a=1/2$ applies: $\rho\Lambda^N\le1/5<1/2$. Explicitly, initialize at mean $x/2$, represent $\Lambda^{-1}R_{2\pi\alpha_i}v_s$ by the vertices for each row, and decode by $2\rho\Lambda^N\ip{e_j}{v_s}$. This verifies every numerical word, including any deterministic remainder of a converse packet decomposition.

There are $(r+1)q$ common command rows and two decoder columns, each row having at most three successors. These counts exclude row-description and arithmetic costs. Since $\eta_\alpha(q)\le1/2$, \eqref{eq:sqrt-upper} follows. Reflection preserves $P_q$; its contracted image uses the same vertex set. The reflected lower bound is obtained simply by restricting to rotation-only words.
\end{proof}

The pair $(\psi_\alpha,\eta_\alpha)$ is a specialization of the word-distortion and enclosure profiles. It is not assumed to characterize exact finite width. In particular, \eqref{eq:psi-lower} can be weak when long word packets have clusters; the optimized profile of Section~\ref{sec:circle} is available without that separation relaxation.

\begin{lemma}[Denominators under a general power bound]\label{lem:weak-denominator}
Suppose for $\tau\ge r$ and $b>0$,
\begin{equation}\label{eq:tau-separation}
        \norm{m\cdot\alpha}_{\R/\Z}\ge b\norm{m}_1^{-\tau}
                         \quad(m\ne0).
\end{equation}
For every integer $Q\ge2$ there is $q$ with
\begin{equation}\label{eq:weak-denominator}
 cQ^{r/(\tau+1-r)}\le q\le Q^r,
 \qquad\eta_\alpha(q)\le Q^{-1},
 \quad c=(b/r^{\tau+1})^{r/(\tau+1-r)}.
\end{equation}
\end{lemma}
\begin{proof}
The $Q^r+1$ points $j\alpha$ in $Q^r$ equal torus boxes give integers $p_i,q$ with $1\le q\le Q^r$ and $|q\alpha_i-p_i|\le Q^{-1}$. For $M=\lfloor q^{1/r}\rfloor$, the $(M+1)^r>q$ integer vectors in $\{0,\ldots,M\}^r$ have a residue collision for $m\cdot p\pmod q$. Subtracting gives $m\ne0$ with $H=\norm{m}_1\le r q^{1/r}$ and $q\mid m\cdot p$. Thus
\[
 bH^{-\tau}\le\norm{m\cdot\alpha}_{\R/\Z}\le H/(qQ),
 \qquad bqQ\le r^{\tau+1}q^{(\tau+1)/r}.
\]
Rearranging yields \eqref{eq:weak-denominator}. This direct residue argument supplies exactly the quantitative transference needed here, without an unspecified transference constant.
\end{proof}

\begin{corollary}[Power bounds at finite Diophantine type]\label{cor:power-type}
Under \eqref{eq:tau-separation}, for constants depending on the fixed data,
\begin{equation}\label{eq:power-bracket}
 c_1N^{r/(2\tau+1)}\le W_{N,\varepsilon}(\alpha)
 \le C_1N^{\min\{1/2,\ r(\tau+1-r)/(\tau+1+r)\}}.
\end{equation}
No matching claim is made when $\tau>r$.
\end{corollary}
\begin{proof}
From \eqref{eq:psi-lower},
$N<B+b^{-2}B^{2\tau+1}\log(1/\kappa)$, while $B\le r(2K)^{1/r}$. This gives the lower power after adjusting finitely many initial horizons. Lemma~\ref{lem:weak-denominator} bounds
$N\eta_\alpha(q)/q^2\le c^{-2}NQ^{-1-2r/(\tau+1-r)}$.
Choose $Q$ at constant times $N^{(\tau+1-r)/(\tau+1+r)}$, large enough to satisfy \eqref{eq:eta-upper} and $q\ge4$. Then $q\le Q^r$ proves the second candidate upper power; \eqref{eq:sqrt-upper} proves the first.
\end{proof}

\begin{corollary}[The retained badly approximable law]\label{cor:ba}
If \eqref{eq:tau-separation} holds with $\tau=r$, then
$W_{N,\varepsilon}(\alpha)=\Theta(N^{r/(2r+1)})$, also with the reflection extension. This is the previously established matched law \cite{GTF42}.
\end{corollary}
\begin{proof}
Substitute $\tau=r$ in \eqref{eq:power-bracket}; the two powers coincide.
\end{proof}

For $1,\alpha_1,\ldots,\alpha_r$ rationally independent, every optimized finite-block action gap on the circle is zero. Indeed simultaneous approximation supplies unbounded frequencies $n$ with every $n\alpha_i$ close to an integer. Each fixed rotation word then acts almost identically on its Fourier mode. For the reflected alphabet use the real cosine mode, which reflection fixes. This is the established almost-invariance argument, not a claim that the finite packet distortion vanishes.
